% Unofficial University of Cambridge Poster Template
% https://github.com/andiac/gemini-cam
% a fork of https://github.com/anishathalye/gemini
% also refer to https://github.com/k4rtik/uchicago-poster
% TeX program = lualatex

\documentclass[final]{beamer}

% ====================
% Packages
% ====================

\usepackage[T1]{fontenc}
\usepackage{lmodern}
\usepackage[orientation=portrait,size=custom,width=120,height=120,scale=1]{beamerposter}
\usetheme{gemini}
\usepackage[dvipsnames]{xcolor}
\usecolortheme{nott}
\usepackage{graphicx}
\usepackage{booktabs}
\usepackage{tikz}
\usetikzlibrary{calc,arrows.meta,patterns,decorations.pathmorphing,shapes.geometric}
\usepackage{tkz-euclide}
\tikzset{point style/.style = {%
  draw = black,
  inner sep = 0pt,
  shape = circle,
  minimum size = 5pt,
  fill = black
 },
 every picture/.append style = {
  scale = 1.5
 },
 every node/.append style={
  scale=1.5
 }
}
\usepackage{pgfplots}
\pgfplotsset{compat=1.14}
\usepackage{anyfontsize}
\usepackage{caption}
\usepackage{subcaption}

% ====================
% Lengths
% ====================

% If you have N columns, choose \sepwidth and \colwidth such that
% (N+1)*\sepwidth + N*\colwidth = \paperwidth
\newlength{\sepwidth}
\newlength{\colwidth}
\setlength{\sepwidth}{0.01\paperwidth}
\setlength{\colwidth}{0.32\paperwidth}

\newcommand{\separatorcolumn}{\begin{column}{\sepwidth}\end{column}}
\newcommand{\bfalert}[1]{\textbf{\alert{#1}}}

% Math shortcuts
\newcommand{\R}{\mathbb{R}}
\newcommand{\N}{\mathbb{N}}
\newcommand{\Z}{\mathbb{Z}}
\newcommand{\Q}{\mathbb{Q}}
\DeclareMathOperator{\s}{succ}
\newcommand{\dv}{\mathbin{\mathrm{div}}}

% Inline shapes
\newcommand{\mysquare}{\tikz[baseline=-7pt]{%
  \node[rectangle,draw,thick,inner sep=6pt] at (0,0) {};
}}
\newcommand{\mytria}{\tikz[baseline=-3.25pt]{%
  \node[isosceles triangle,isosceles triangle apex angle=60,draw,thick,inner
  sep=3.25pt,rotate=90] at (0,0) {};
}}
\newcommand{\mycirc}{\tikz[baseline=-7pt]{%
  \node[circle,draw,thick,inner sep=4.5pt,baseline=0.5ex,rotate=90] at (0,0) {};
}}
\newcommand{\mycross}{\tikz[baseline=-7pt,scale=0.2]{%
  \draw[thick] (-1,1) -- (1,-1);
  \draw[thick] (-1,-1) -- (1,1);
}}

% Colors %
\newcommand{\clr}{\textcolor{BrickRed}}
\newcommand{\clb}{\textcolor{RoyalBlue}}
\newcommand{\clg}{\textcolor{ForestGreen}}
\newcommand{\clm}{\textcolor{Fuchsia}}
\newcommand{\clv}{\textcolor{violet}}
\newcommand{\clbr}{\textcolor{Sepia}}
\newcommand{\cly}{\textcolor{Dandelion}}

% ====================
% Title
% ====================

\title{Number Theory Cheatsheet}

\author{3.AB PreIB Math}

\institute[shortinst]{Adam Klepáč and Jáchym Löwenhöffer}

% ====================
% Footer (optional)
% ====================

% \footercontent{
%   \href{https://utfpr.edu.br/ct/ppgca}{utfpr.edu.br/ct/ppgca} \hfill
%   Mostra de Trabalhos do PPGCA --- TechTalks 2024 \hfill
%   \href{mailto:ppgca-ct@utfpr.edu.br}{ppgca-ct@utfpr.edu.br}}
% (can be left out to remove footer)


% ====================
% Logo (optional)
% ====================

% use this to include logos on the left and/or right side of the header:
\logoright{\includegraphics[height=6cm]{logos/logo-white.png}}
% \logoleft{\hspace{20ex}\includegraphics[height=3.5cm]{logos/ppgca-logo.png}}

% ====================
% Body
% ====================

\begin{document}

% Refer to https://github.com/k4rtik/uchicago-poster
% logo: https://www.cam.ac.uk/brand-resources/about-the-logo/logo-downloads
% \addtobeamertemplate{headline}{}
% {
%     \begin{tikzpicture}[remember picture,overlay]
%       \node [anchor=north west, inner sep=3cm] at ([xshift=-2.5cm,yshift=1.75cm]current page.north west)
%       {\includegraphics[height=7cm]{logos/unott-logo.eps}}; 
%     \end{tikzpicture}
% }

\begin{frame}[t]
\begin{columns}[t]
\separatorcolumn

\begin{column}{\colwidth}

 \begin{exampleblock}{Natural Numbers}
  Natural numbers (denoted $\N$) are defined basically as `sets containing so
  many elements'. This means that the number $0$ is a set with no elements, $1$
  is a set with one element and so on. Formally, we construct them in the
  following way ($\emptyset$ is the empty set):
  \[
   \begin{array}{r c l c l}
    0 & = & \emptyset & &\\
    1 & = & \{0\} & = & \{\emptyset\}\\
    2 & = & \{0, 1\} & = & \{\emptyset, \{\emptyset\}\}\\
    3 & = & \{0, 1, 2\} & = & \{\emptyset, \{\emptyset\}, \{\emptyset,
    \{\emptyset\}\}\}\\
      & \vdots & & & 
   \end{array}
  \]
  In general, the \alert{next natural number} after a number $n$ is defined as
  the set $\{0,\ldots,n\}$. 

  Observe that we can find a formula for the next number after $n$. Since $n =
  \{0,\ldots,n-1\}$ and the next number is $\{0,\ldots,n\}$, we can construct
  the next number after $n$ as a union of two sets: $n \cup \{n\}$. We call this
  number, the \alert{successor} of $n$, and write it as \alert{$\s(n)$}. For
  example, $1 = \{0\} = 0 \cup \{0\} = \s(0)$ or $3 = \{0,1,2\} = \{0,1\} \cup
  \{2\} = 2 \cup \{2\} = \s(2)$.

  \vspace{18pt}

  \textbf{\large Addition of natural numbers (\alert{not examined})}

  We can define the operation of \alert{addition} on natural numbers using two
  simple rules. For two natural numbers $n,m \in \N$,
  \begin{enumerate}[label=(\arabic*),left=12pt]
   \item $n + 1 = \s(n)$,
   \item $\s(n + m) = n + \s(m)$.
  \end{enumerate}
  Rule (1) simply states that $n + 1$ is the next number after $n$. Rule (2) is
  harder to decode. It literally says that by adding the two numbers together
  and then taking the next number one reaches the same answer as by first taking
  the next number and then performing addition. It's visualised on the picture
  below.
  \begin{figure}[H]
   \centering
   \begin{tikzpicture}
    \begin{scope}
     \foreach \x in {-.5, 0, .5} {
      \node[circle,fill=BrickRed,inner sep=3pt] at (\x,.75) {};
     }
     \node[BrickRed] at (0,0) {\footnotesize $n$};

     \foreach \x in {1.5, 2, 2.5, 3} {
      \node[circle,fill=RoyalBlue,inner sep=3pt] at (\x,.75) {};
     }
     \node[RoyalBlue] at (2.25,0) {\footnotesize $m$};
    \end{scope}
    \begin{scope}[xshift=5.5cm]
     \foreach \x in {-.5, 0, .5} {
      \node[circle,fill=BrickRed,inner sep=3pt] at (\x,.75) {};
     }
     \node at (.25,0) {\footnotesize $\clr{n} + \clb{m}$};

     \foreach \x in {-.5, 0, .5, 1} {
      \node[circle,fill=RoyalBlue,inner sep=3pt] at (\x,1.25) {};
     }
    \end{scope}
    \begin{scope}[xshift=8.5cm]
     \foreach \x in {-.5, 0, .5} {
      \node[circle,fill=BrickRed,inner sep=3pt] at (\x,.75) {};
     }
     \node at (.25,0) {\footnotesize $\clg{\s}(\clr{n} + \clb{m})$};

     \foreach \x in {-.5, 0, .5, 1} {
      \node[circle,fill=RoyalBlue,inner sep=3pt] at (\x,1.25) {};
     }
     \node[circle,fill=ForestGreen,inner sep=3pt] at (-.5, 1.75) {};
    \end{scope}
    \begin{scope}[xshift=12.5cm]
     \foreach \x in {-.5, 0, .5, 1} {
      \node[circle,fill=RoyalBlue,inner sep=3pt] at (\x,.75) {};
     }
     \node[circle,fill=ForestGreen,inner sep=3pt] at (-.5, 1.25) {};
     \node at (.25,0) {\footnotesize $\clg{\s}(\clb{m})$};
    \end{scope}
    \begin{scope}[xshift=15.5cm]
     \foreach \x in {-.5, 0, .5} {
      \node[circle,fill=BrickRed,inner sep=3pt] at (\x,.75) {};
     }
     \node[circle,fill=ForestGreen,inner sep=3pt] at (-.5, 1.75) {};
     \node at (.25,0) {\footnotesize $\clr{n} + \clg{\s}(\clb{m})$};

     \foreach \x in {-.5, 0, .5, 1} {
      \node[circle,fill=RoyalBlue,inner sep=3pt] at (\x,1.25) {};
     }
    \end{scope}
   \end{tikzpicture}

   \caption{Visualisation of rule (2) of addition. Both $\clg{\s}(\clr{n} +
    \clb{m})$ and $\clr{n} + \clg{\s}(\clb{m})$ feature the \alert{same number}
    of dots.}
   \label{fig:rule-2-addition}
  \end{figure}

  Rules (1) and (2) combine to give a simple algorithm of computing the sum $n +
  m$ for any two numbers $n,m \in \N$. It goes like this:
  \begin{itemize}[label=\textbullet,left=12pt]
   \item Using rule (1), calculate $n + 1 = \s(n) = n \cup \{n\}$.
   \item Now that we have calculated $n + 1$, we can calculate $n + 2$ because
    $n + 2 = n + \s(1)$ and by rule (2) this equals $n + \s(1) = \s(n + 1)$, so
    just take the next number after $n + 1$.
   \item Continue like this until you calculate $n + m = n + \s(m - 1)$.
  \end{itemize}
  For example, to compute $4 + 2$, we calculate $4 + 1 = \s(4)$ and then $4 + 2
  = 4 + \s(1) = \s(4 + 1) = \s(\s(4))$ so $4 + 2$ is just the next number after
  the next number after $4$.
 \end{exampleblock}

 \begin{exampleblock}{Integers (Whole Numbers)}
  We have defined \alert{addition} on natural numbers, but in order to perform
  \alert{subtraction}, we must move to a `larger' set of numbers -- the
  \alert{integers}. This is because subtraction is \alert{\textbf{not}} an
  operation on natural numbers as its result needn't be a natural number itself.

  The idea behind the definition of integers (labelled $\Z$) is to take
  \alert{pairs of natural numbers}. Fundamentally, we want the pair $(a,b) \in
  \N \times \N$ to \alert{represent} the result of the operation `$a-b$' (which
  we can't yet perform because we need to define the integers \textbf{before}
  defining subtraction).

  To this end, we define an \alert{equivalence} on $\N \times \N$ (i.e. on pairs
  of natural numbers) that makes two pairs equivalent \alert{if they represent
  the same integer}. For example, the pair $(4, 6)$ should represent the number
  $-2$ (as $4 - 6 = -2$) and so should the pairs $(8, 10)$, $(3, 5)$ or just
  about any pair $(a, a + 2)$ for $a \in \N$. The integers will then be the
  \alert{classes of equivalence} of this equivalence relation.

  We label this equivalence by the letter $\clg{E}$. Since we want $(a,b)$ to be
  \clg{equivalent} to $(c,d)$ if `$a - b = c - d$' but we can't use subtraction
  yet, we simply rewrite the equation above to use only addition, like this: $a
  + d = c + b$. Thus, we say that $\clr{(a,b)}\clg{E}\clb{(c,d)}$ if $\clr{a} +
  \clb{d} = \clb{c} + \clr{b}$. This defines an equivalence on $\N \times \N$
  and we let $\Z$ be the classes of equivalence of all pairs of natural numbers:
  \[
   \clm{\Z} = \{[(a,b)]_{\clg{E}} \mid a,b \in \N\}.
  \]
  To give an example, the pair $\clr{(3,5)}$ is \clg{equivalent} to
  $\clb{(7,9)}$ because $\clr{3} + \clb{9} = \clb{7} + \clr{5}$ and they both
  represent the integer $\clm{-2}$. Similarly, both $\clr{(6, 2)}$ and
  $\clb{(8, 4)}$ represent the integer $\clm{4}$. The visualisation of integers
  as pairs of natural numbers is given below.
  \begin{figure}[H]
   \centering
   \begin{tikzpicture}
    \node[ForestGreen] at (-1, 2.5) {\clg{$\cdots$}};
    \draw[ForestGreen,thick] (0, 5) -- (0, 0);
    \begin{scope}[xshift=1.5cm,yshift=1cm]
     \node[BrickRed] at (0, 3) {\footnotesize $(2, 4)$};
     \node[BrickRed] at (2, 3) {\footnotesize $(7, 9)$};
     \node[BrickRed] at (0, 2) {\footnotesize $(0, 2)$};
     \node[BrickRed] at (2, 2) {\footnotesize $(11, 13)$};
     \node at (1, 1) {\footnotesize $\vdots$};
     \node at (1, 0) {\footnotesize all pairs $\clr{(a, a + 2)}$};
     \node at (1, -2) {$\clm{-2}$};
    \end{scope}
    \draw[ForestGreen,thick] (5, 5) -- (5, 0);
    \begin{scope}[xshift=6.5cm,yshift=1cm]
     \node[BrickRed] at (0, 3) {\footnotesize $(3, 4)$};
     \node[BrickRed] at (2, 3) {\footnotesize $(10, 11)$};
     \node[BrickRed] at (0, 2) {\footnotesize $(0, 1)$};
     \node[BrickRed] at (2, 2) {\footnotesize $(22, 23)$};
     \node at (1, 1) {\footnotesize $\vdots$};
     \node at (1, 0) {\footnotesize all pairs $\clr{(a, a + 1)}$};
     \node at (1, -2) {$\clm{-1}$};
    \end{scope}
    \draw[ForestGreen,thick] (10, 5) -- (10, 0);
    \node[ForestGreen] at (11, 2.5) {\clg{$\cdots$}};
    \draw[ForestGreen,thick] (12, 5) -- (12, 0);
    \begin{scope}[xshift=13.5cm,yshift=1cm]
     \node[BrickRed] at (0, 3) {\footnotesize $(8, 5)$};
     \node[BrickRed] at (2, 3) {\footnotesize $(13, 10)$};
     \node[BrickRed] at (0, 2) {\footnotesize $(3, 0)$};
     \node[BrickRed] at (2, 2) {\footnotesize $(7, 4)$};
     \node at (1, 1) {\footnotesize $\vdots$};
     \node at (1, 0) {\footnotesize all pairs $\clr{(a + 3, a)}$};
     \node at (1, -2) {$\clm{3}$};
    \end{scope}
    \draw[ForestGreen,thick] (17, 5) -- (17, 0);
    \node[ForestGreen] at (18, 2.5) {\clg{$\cdots$}};
   \end{tikzpicture}
   \caption{\clm{Integers} as classes of \clg{equivalence} of natural numbers.}
   \label{fig:integers}
  \end{figure}
  The \alert{addition} of integers is defined using the addition of natural
  numbers. Given two classes of equivalence $[(\clr{a},\clr{b})]_{\clg{E}},
  [(\clb{c},\clb{d})]_{\clg{E}} \in \clm{\Z}$, we let
  \[
   [(\clr{a},\clr{b})]_{\clg{E}} + [(\clb{c}, \clb{d})]_{\clg{E}} = [(\clr{a} +
   \clb{c}, \clr{b} + \clb{d})]_{\clg{E}}.
  \]
  Finally, we define the \alert{opposite number} to $[(a,b)]_{\clg{E}}$ as
  $-[(a,b)]_{\clg{E}} = [(b,a)]_{\clg{E}}$ (this is because $-(a-b) = b-a$). The
  \alert{subtraction} of two integers is now just a sum of the first and the
  opposite of the second, that is
  \[
   [(\clr{a},\clr{b})]_{\clg{E}} - [(\clb{c},\clb{d})]_{\clg{E}} =
   [(\clr{a},\clr{b})]_{\clg{E}} + (-[(\clb{c},\clb{d})]_{\clg{E}}) =
   [(\clr{a},\clr{b})]_{\clg{E}} + [(\clb{d},\clb{c})]_{\clg{E}} = [(\clr{a} +
   \clb{d}, \clr{b} + \clb{c})]_{\clg{E}}.
  \]
  For example,
  \[
   [(\clr{3}, \clr{1})]_{\clg{E}} - [(\clb{5},\clb{2})]_{\clg{E}} = [(\clr{3},
   \clr{1})]_{\clg{E}} + [(\clb{2},\clb{5})]_{\clg{E}} = [(\clr{3} + \clb{2},
   \clr{1} + \clb{5})]_{\clg{E}} = [(5, 6)]_{\clg{E}},
  \]
  which is the same as writing
  \[
   \clr{2} - \clb{3} = \clr{2} + (-\clb{3}) = -1.
  \]
 \end{exampleblock}

\end{column}

\separatorcolumn

\begin{column}{\colwidth}

\begin{block}{Multiplication}
 In a way similar to addition, we can define \alert{multiplication} on natural
 numbers by the following two rules.
 \begin{enumerate}[label=(\arabic*),left=12pt]
  \item $n \cdot 1 = n$,
  \item $n \cdot \s(m) = n \cdot m + m$,
 \end{enumerate}
 for $n,m \in \N$. They carry the idea behind an algorithmic way to compute the
 product $n \cdot m$ for any two natural numbers $n,m$. It goes like this:
 \begin{itemize}[left=12pt,label=\textbullet]
  \item Using rule (1), calculate $n \cdot 1 = n$.
  \item Using rule (2), calculate $n \cdot 2 = n \cdot \s(1) = n \cdot 1 + n = n
   + n$.
  \item Continue like this until you calculate
   \[
    n \cdot m = n \cdot \s(m - 1) = n \cdot (m - 1) + n = \underbrace{n + n +
    \ldots + n}_{(m - 1)\text{ times}} + n.
   \]
 \end{itemize}
 For example, to calculate $4 \cdot 3$, we first multiply $4 \cdot 1 = 4$, then
 $4 \cdot 2 = 4 \cdot \s(1) = 4 \cdot 1 + 4 = 4 + 4$, and finally $4 \cdot 3 = 4
 \cdot \s(2) = 4 \cdot 2 + 4 = (4 + 4) + 4$. As you've been taught:
 `multiplication is just repeated addition'.

 Multiplication is easily extended to integers by the formula
 \[
  [(\clr{a},\clr{b})]_{\clg{E}} \cdot [(\clb{c},\clb{d})]_{\clg{E}} = [(\clr{a}
  \cdot \clb{c} + \clr{b} \cdot \clb{d}, \clr{b} \cdot \clb{c} + \clr{a} \cdot
  \clb{d})]_{\clg{E}}.
 \]
 The formula is based on the calculation
 \[
  (\clr{a} - \clr{b}) \cdot (\clb{c} - \clb{d}) = \clr{a} \cdot \clb{c} -
  \clr{b} \cdot \clb{c} - \clr{a} \cdot \clb{d} + \clr{b} \cdot \clb{d} =
  (\clr{a} \cdot \clb{c} + \clr{b} \cdot \clb{d}) - (\clr{b} \cdot \clb{c} +
  \clr{a} \cdot \clb{d}).
 \]

 For example,
 \[
  [(\clr{5},\clr{3})]_{\clg{E}} \cdot [(\clb{1},\clb{5})]_{\clg{E}} =
  [(\clr{5} \cdot \clb{1} + \clr{3} \cdot \clb{5}, \clr{3} \cdot \clb{1} +
  \clr{5} \cdot \clb{5})]_{\clg{E}} = [(20, 28)]_{\clg{E}}.
 \]
 This is the same calculation as
 \[
  \clr{2} \cdot (\clb{-4}) = -8.
 \]
\end{block}

\begin{exampleblock}{Rational Numbers}
 Being able to \alert{multiply integers}, we'd like to divide them as well. As
 was the case with natural numbers and subtraction, \alert{division is not an
 operation on integers} because its result needn't be an integer.

 The idea behind the definition of rational numbers (labelled $\Q$) is pretty
 much the same as the one behind the definition of integers -- rational numbers
 are really just \alert{pairs of integers}. And again, multiple pairs of
 integers \alert{represent the same} rational number. Therefore, given pairs
 $(a,b)$ and $(c,d)$ with $a,b,c,d \in \Z$, we must make sure that $(a,b)$
 \alert{is equivalent to} $(c,d)$ if `the fraction $a / b$ is the same as the
 fraction $c / d$'.

 As we couldn't have defined division yet, we must \alert{rewrite the last
 equation in terms of multiplication only}. This is easy to do because $a / b =
 c / d$ if $a \cdot d = c \cdot b$. This directly leads to the definition of an
 \alert{equivalence} $\clg{Q}$ on pairs of integers: $(\clr{a},
 \clr{b})\clg{Q}(\clb{c},\clb{d})$ if
 \[
  \clr{a} \cdot \clb{d} = \clr{c} \cdot \clb{b}.
 \]
 This is indeed an equivalence on $\Z \times \Z$ and we define $\clm{\Q}$ as
 \[
  \clm{\Q} = \{[(a,b)]_{\clg{Q}} \mid a,b \in \Z\}.
 \]
 We tend to write elements of $\clm{\Q}$ as \textbf{fractions}, that is, instead
 of $[(a,b)]_{\clg{Q}}$, we write $a / b$. We shall adopt this notation
 henceforth.

 It only remains to \alert{extend addition and multiplication} to rational
 numbers. This is easily done using formulae you already know. For example, the
 \alert{product} of two rational numbers $a / b, c / d \in \clm{\Q}$ is defined
 as such:
 \[
  \frac{\clr{a}}{\clr{b}} \cdot \frac{\clb{c}}{\clb{d}} = \frac{\clr{a} \cdot
  \clb{c}}{\clr{b} \cdot \clb{d}}.
 \]
 The \alert{sum} of rational numbers as
 \[
  \frac{\clr{a}}{\clr{b}} + \frac{\clb{c}}{\clb{d}} = \frac{\clr{a} \cdot
  \clb{d} + \clb{c} \cdot \clr{b}}{\clr{b} \cdot \clb{d}}.
 \]
 For example,
 \[
  \frac{\clr{2}}{\clr{5}} \cdot \frac{\clb{3}}{\clb{4}} = \frac{\clr{2} \cdot
  \clb{3}}{\clr{5} \cdot \clb{4}} = \frac{6}{20} \quad \text{and} \quad
  \frac{\clr{2}}{\clr{5}} + \frac{\clb{3}}{\clb{4}} = \frac{\clr{2} \cdot
  \clb{4} + \clb{3} \cdot \clr{5}}{\clr{5} \cdot \clb{4}} = \frac{23}{20}.
 \]
 
 Finally, we're ready to \alert{define division} on rational numbers. We first
 define the \alert{inverse} of a rational numbers $a / b$ as $b / a$. We write
 $b / a = (a / b)^{-1}$. The \alert{operation of division} on rational numbers
 is defined as \alert{multiplication by the inverse element}, that is
 \[
  \frac{\clr{a}}{\clr{b}} : \frac{\clb{c}}{\clb{d}} = \frac{\clr{a}}{\clr{b}}
  \cdot \left( \frac{\clb{c}}{\clb{d}} \right)^{-1} = \frac{\clr{a}}{\clr{b}}
  \cdot \frac{\clb{d}}{\clb{c}} = \frac{\clr{a} \cdot \clb{d}}{\clr{b} \cdot
  \clb{c}}.
 \]
 For example,
 \[
  \frac{\clr{2}}{\clr{5}} : \frac{\clb{3}}{\clb{4}} = \frac{\clr{2}}{\clr{5}}
 \cdot \left(\frac{\clb{3}}{\clb{4}}\right)^{-1} = \frac{\clr{2}}{\clr{5}} \cdot
 \frac{\clb{4}}{\clb{3}} = \frac{\clr{2} \cdot \clb{4}}{\clr{5} \cdot \clb{3}} =
 \frac{8}{15}.
 \]
\end{exampleblock}

\begin{alertblock}{Divisibility}
 There are two more interesting operations on integers -- \alert{integer
 division} and \alert{modulus}. You've learnt about them in elementary school
 and they are basically `a way to do \alert{division on integers}'.

 Given two integers: $a,b \in \Z$, we can ask `\textbf{How many times does $b$
 fit into $a$?}'. This number is called the \alert{quotient} (of $a$ by $b$) and
 denoted \alert{$a \dv b$}. For example, $7 \dv 2 = 3$ because $2$ fits $3$
 times into $7$ or $10 \dv 8 = 1$ because $8$ fits only once into $10$. This
 operation can also be used with integers, for instance $-5 \dv 2 = -2$ because
 $2$ `fits' $-2$ times into $-5$.

 The quantity that is `left after integer division' is called the
 \alert{remainder} and denoted $a \bmod b$. In our first example, $7 \bmod 2 =
 1$ because $7 \dv 3 = 2$ and $2 \cdot 3 = 6$, so the number $1$ is left after
 performing the division. Similarly, $10 \bmod 8 = 2$ since $10 \dv 8 = 1$ and
 $8 \cdot 1 = 8$ and $10 - 8 = 2$. In the last example, we get $-5 \bmod 2 = -1$
 as $-5 \dv 2 = -2$ and $2 \cdot (-2) = -4$.

 Notice that the \alert{remainder must always be smaller (in absolute value)
 than the divisor} because it's the quantity that's left after the divisor no
 longer fits into the dividend. We may thus formalise the idea of \alert{integer
 division as such}: for any two numbers $a,b \in \Z$ there are \textbf{unique}
 numbers $q,r \in \Z$ (called \alert{quotient} and \alert{remainder}) such that
 \[
  a = b \cdot q + r
 \]
 and $0 \leq |r| < |b|$. We write $q = a \dv b$ and $r = a \bmod b$.

 The operation of integer division gives birth to the idea of
 \alert{divisibility}. We say that a number $b$ divides the number $a$ (and
 write $b \mid a$) if $a \bmod b = 0$ or, equivalently, $a = b \cdot q$ for some
 integer $q \in \Z$.
\end{alertblock}

\end{column}
\separatorcolumn

\begin{column}{\colwidth}

\begin{exampleblock}{Greatest Common Divisor}
 We ask the question: `What's the \alert{largest number that divides} both $a$
 and $b$?' Such number is called the \alert{greatest common divisor} of $a$ and
 $b$ and written $\mathrm{gcd}(a,b)$.

 There is a nice \alert{algorithmic way} to compute this number: the algorithm
 is called \alert{Euclid's Algorithm}. It uses the following equality: for every
 two numbers, $a,b \in \Z$, it holds that
 \[
  \mathrm{gcd}(a,b) = \mathrm{gcd}(a \bmod b, b).
 \]
 Since $a \bmod b$ is always smaller than both $a$ and $b$, we keep taking the
 remainder after division of one number by the other until we reach the number
 $0$. Once we do, we have found the greatest common divisor as
 $\mathrm{gcd}(d,0) = d$ ($0$ is divided by every number).

 Let's showcase the algorithm. Suppose we want to calculate $\mathrm{gcd}(3312,
 448)$. We first calculate $3312 \bmod 448 = 176$. Therefore
 \[
  \mathrm{gcd}(3312, 448) = \mathrm{gcd}(176, 448).
 \]
 Next, we compute $448 \bmod 176 = 96$ and thus
 \[
  \mathrm{gcd}(176, 448) = \mathrm{gcd}(176, 96).
 \]
 We are almost done. Computing $176 \bmod 96 = 80$ and $96 \bmod 80 = 16$ gives
 \[
  \mathrm{gcd}(176, 96) = \mathrm{gcd}(80, 96) = \mathrm{gcd}(80, 16).
 \]
 Because $80 \bmod 16 = 0$, we have reached the conclusion that
 $\mathrm{gcd}(3312, 448) = 16$.
\end{exampleblock}

\end{column}
\separatorcolumn

\end{columns}
\end{frame}

\end{document}
